
\documentclass[class=llncs, crop=false]{standalone}
\input{preamble.tex}
\input{macros.tex}

\begin{document}
%
\section{Encoding Eunoia in the λΠ-calculus}\label{sec:lambdapi}
\begin{boxfigure}[t!]{fig:lpTermSyntax}
	{Syntax and typing rules for the $λΠ$-calculus modulo rewriting.}
	\alttext{Top: term grammar for the lambda-Pi calculus modulo rewriting.
		Universes mu range over TYPE and KIND.
		Terms t range over variables, constants, universes, applications,
		lambda abstractions, and dependent products.
		Bottom: natural-deduction typing rules---empty context is well-formed
		(wf0); extending with (x:t) requires Gamma deriving t at some universe
		(wf+); variable rule reads (x:t) from the context (var);
		constant rule reads (kappa:t) from the signature with the constant's
		type sortable in Gamma (con); TYPE has type KIND (univ);
		dependent product, abstraction, application, and conversion modulo
		beta and rewriting (prod, fun, app, conv).}{%
		\vspace{4pt}\centerline{$\lpTermSyntax$}
		\boxfigurehline
		\centerline{$\lpTypingRules$}\vspace{4pt}}
\end{boxfigure}
%
\Autoref{fig:lpTermSyntax} provides syntax and typing rules
for the \emph{$λΠ$-calculus modulo rewriting}.
%
Each term is either a \emph{variable} $x$,
a \emph{constant} $κ$,
a \emph{universe} from $\{ \TYPE, \KIND \}$,
an \emph{application} of two terms $\prn{\app{t}{t'}}$,
or an \emph{abstraction}
$\prn{\binddot{\mscr{B}}{x : t}{t'}}$
where $\mscr{B}$ is a \emph{binder} from $\{ λ, Π \}$.
%
Terms are identified up to $α$-conversion
(i.e., renaming of bound variables), and we assume
appropriate definitions for \emph{substitution} $\prn{t[x ↦ t']}$
and \emph{$β$-reduction} $\prn{t ⇝_{β} t'}$.
%
Hereinafter, we may write $\prn{t_1 → t_2}$ for the term
$\prn{\ppii{(x : t_1)}{t_2}}$, where $x$ does not occur
free in $t_2$.

% --------------------------------------------------------%
A \emph{typing} is an expression of the form $\prn{t : t'}$,
and a \emph{context} is a list of typings each of the
form $\prn{x : t}$.
%
% Hereinafter, given a context $Γ$ and a typing $\prn{x : t}$,
% let $\prn{Γ, \prn{x : t}}$ denote the extension $Γ ∪ \{x : t\}$.
%
A \emph{rewrite rule} is an expression of the form
$\prn{t ↪ t'}$, where $t$ has the form
$\prn{\app{\prn{\app{κ}{t_1}}\ \ldots}{t_n}}$.
%
A \emph{signature} is a list of typings
$\prn{κ : t}$ (with $t$ closed) and rewrite rules.
%
For any signature $Σ$, let $R_{Σ}$ be the smallest
binary relation such that:
%
\begin{enumerate}
	\item for any rule $\prn{ℓ ↪ r} ∈ Σ$ and substitution $σ$,
	      $(ℓσ, rσ) ∈ R_{Σ}$, and
	\item $R_{Σ}$ is \emph{congruent} under application
	      and abstraction.
\end{enumerate}
%
%, abstraction and substitution.  $R$
Then, \emph{equality modulo rewriting} $\prn{{≡_{βΣ}}}$ is
defined as the least equivalence relation containing
$R_{Σ}$ and the $β$-reduction relation.
%
The rules in \autoref{fig:lpTermSyntax}
provide a (mutually inductive) definition of a
\emph{well-formedness} relation $\prn{\msf{wf}_Σ}$
on contexts and a \emph{typing relation}
$\prn{⊢_{Σ}}$ between contexts and typings,
where $\prn{Γ ⊢_{Σ} e : t}$ may be read as
``$e$ has type $t$ with respect to signature $Σ$ and context $Γ$''.

% -------------------------------------------------------- %
\subsection{Translating Eunoia Terms}
%
The typing rules of the $λΠ$-calculus disallow
abstractions on types (i.e., there are no expressions
of type $\TYPE → t$ in any context). Eunoia's types are
therefore encoded as $λΠ$-terms of type $\Set$, using
the Tarski-style universe of LambdaPi's standard
library: an injective symbol
$\msf{τ} : \Set → \TYPE$ lifts a $\Set$-term to a
$λΠ$-type, and the (infix) function-space symbol
$\prn{⤳} : \Set → \Set → \Set$ satisfies
$\msf{τ}\,\prn{A ⤳ B} ↪ \msf{τ}\,A → \msf{τ}\,B$.
%
The dependent types arising from $\eo{quote}$ arguments
are not expressible with $\prn{⤳}$ alone, so we
additionally introduce a dependent function space
$\prn{⤳^{\msf d}} :
	\pii{T : \Set}{\prn{{\msf{τ}\,T → \Set}} → \Set}$
with the rule
$\msf{τ}\,\prn{T ⤳^{\msf d} F} ↪
	\pii{x : \msf{τ}\,T}{\msf{τ}\,\prn{\app F x}}$;
the non-dependent case is recovered by taking
$F ≔ \lam{\_}{B}$.
%
We keep $\prn{⤳}$ as a first-class operator rather than
defining it from $\prn{⤳^{\msf d}}$ because programs
need a matchable function-space head: pattern variables
cannot bind under a $Π$, but they can under $\prn{⤳}$.

The translation $\tm{⋅}$ on Eunoia terms below produces
$λΠ$-terms in the abstract calculus of
\autoref{fig:lpTermSyntax} and assumes all terms have
already been processed by the $\mbf{elab}$ operator of
\autoref{sec:eo-elab}.
%
Because some Eunoia symbols are not valid LambdaPi
identifiers, the translation on symbols uses LambdaPi
syntax for escaping: $\bar s ≔ \esc s$ iff $s$
contains forbidden characters, and $\bar s ≔ s$
otherwise.
%
\begin{definition}\label{def:tm-ty}
	%
	Let $\tm{⋅}$ be the least operator such that:

	\begin{enumerate}
		\item For any symbol $s$ or literal $ℓ$, the following hold:
		      $$\tm{s} =
			      \begin{scases}
				      \lpo{Type}   & \text{if $s = \eotype$,} \\
				      \bar s & \text{otherwise.}
			      \end{scases}
			      %
			      \quad
			      %
			      \tm{ℓ} =
			      \begin{scases}
				      \ev{of\_Z}\ \tm{n}
				      &
				      \text{if $ℓ = \msf{int}(n)$,}
				      \\
				      \ev{of\_Q}\ \tm{n}\ \tm{m}
				      &
				      \text{if $ℓ = \msf{rat}(n,m)$,}
				      \\
				      \estr{s} & \text{if $ℓ = \msf{str}(s)$.}
			      \end{scases}
		      $$
		      The symbols $\ev{of\_Z}$ and $\ev{of\_Q}$ are
		      defined in our prelude (\autoref{sec:results}) and embed
		      the integer and rational literals of LambdaPi's
		      standard library into the carrier type of Eunoia's
		      arithmetic sorts.
		\item For arrow types, the following hold,
		      where $\epar{x\ T} ∈ \vec ρ$:
		      $$
			      \begin{array}[t]{r@{\ {=}\ }l}
				      \tm{\earr{\eoquote{x}\ {\vec t}}}
				                    &
				      \tm{T} ⤳^{\msf d} \lam{x}{\tm{\earr{\vec t}}}
				      \\[1mm]
				      \tm{\earr{t'\ {\vec t}}}
				                    &
				      \tm{t'} ⤳ \tm{\earr{\vec t}}
				      \\[1mm]
				      \tm{\earr{t}} & \tm t
			      \end{array}
		      $$
		\item For binary and $n$-ary applications,
		      the following hold:
		      $$
			      \begin{array}[t]{l@{\ {=}\ }l}
				      \tm{\eapp{\mtt{\_}}{t_1\ t_2}}
				       & \lpo{app}\ \tm{t_1}\ \tm{t_2}
				      \\[2mm]
				      \tm{\eprf{t}}
				       & \lpo{Proof}\ {\tm {t}}
				      \\[2mm]
				      \tm{\eapp{s}{\vec t}}
				       &
				      \bar s\ \tm{t_1}\ \ldots\ \tm{t_n}
			      \end{array}
		      $$
		\item After elaboration, every other binder
		      application is unfolded by item~2 of
		      \autoref{def:elab}; only $\eo{define}$
		      survives because it is a let-binder rather
		      than a binding form. It is translated using
		      LambdaPi's standard $\msf{let}$ binder:
		      $$\tm{\eapp{\eo{define}}{
					      \epar{\vec v}
					      \ t'
				      }
			      }
			      = \prn{
				      \msf{let}\ \vec {w}
				      \ \msf{in}\ \tm{t'}
			      }$$
		      where $w_i = \prn{{\bar x}_i ≔ \tm{t_i}}$
		      for each $v_i = \epar{x_i\ t_i}$.
	\end{enumerate}
	%
\end{definition}

% -------------------------------------------------------- %
\subsection{The LambdaPi Proof Assistant}
%
\begin{boxfigure}[t!]{fig:lpSyntax}
	{Syntax for a fragment of the LambdaPi proof assistant.}
	\alttext{Two grammar tables for LambdaPi concrete syntax.
		Top: modifiers (constant, sequential); parameters (explicit (s:t) and
		implicit [s:t]); patterns theta (pattern variables \$x or pattern
		applications); rewrite rules (s applied to patterns rewriting to theta).
		Bottom: commands---symbol declaration with optional modifier, parameters,
		type, and optional definition; rule declaration listing rewrite rules with
		'with' separators; assertion 'assert |- t : t';
		and 'require open' for module imports.}{%
		\vspace{4pt}\centerline{$\lpPatternSyntax$}
		\boxfigurehline
		\centerline{$\lpConcreteSyntax$}\vspace{4pt}}
\end{boxfigure}
%
\Autoref{fig:lpSyntax} defines the fragment of LambdaPi
required for encoding Eunoia. It is sugar over the
calculus of \autoref{fig:lpTermSyntax}: implicit binders
$\imp{x : t}$ and placeholders $\_$ are elaborated to
$Π$-binders and metavariables during type inference, and
each declaration extends~$Σ$.

A \emph{symbol declaration}
$\maybe{m}\,\kw{symbol}\ s\ \vec ρ : t\,\maybe{≔ t'}$
introduces $s$ with the typing
$\prn{\xsym s : \pii{\vec ρ}{t}}$, where parameters
$\vec ρ$ are either \emph{explicit} $\xpl{x : t}$ or
\emph{implicit} $\imp{x : t}$.
We write $\xsym s$ (with an @) for the
\emph{fully-explicit} form; an unprefixed~$s$ inserts a
placeholder for each leading implicit, to be filled by
higher-order unification.
The optional $\md{constant}$ modifier forbids rewrite
rules with $s$ at the head; $\md{sequential}$ matches
rewrite rules with head $s$ in their textual order; and
a definition $\prn{≔ t'}$ adds the rewrite rule
$\prn{s ↪ \lam{\vec ρ}{t'}}$.
%
A \emph{rewrite rule declaration} $\kw{rule}\ θ ↪ θ'$
adds the rule $\prn{θ ↪ θ'}$ to the signature, where
$θ$ and $θ'$ are \emph{patterns}---either a
\emph{pattern variable} $\ptrn{x}$ (written \$x) or a
\emph{pattern application} $\prn{s\,\any{θ}}$.
The rule is valid when $θ$ is a pattern application and
every pattern variable in $θ'$ also occurs in $θ$.
LambdaPi additionally verifies \emph{subject reduction}
at declaration time: for every typed substitution of
the pattern variables, both sides must have the same
type. An \emph{assertion} $\kw{assert}\ ⊢\ t : t'$
checks $t : t'$ without extending the signature.

% ----------------------------------------------------------
\subsection{Translating Eunoia Commands and Proofs}
% ----------------------------------------------------------
We extend the term translation $\tm{⋅}$
(\autoref{def:tm-ty}) to parameter contexts and to
both standard and proof commands.
%
\Autoref{tab:encoding} summarises the full
Eunoia-to-LambdaPi encoding; the remainder of this
section formalises the operators on contexts
(\autoref{def:ctx}), standard commands
(\autoref{def:cmd}), and proof commands
(\autoref{def:cmd-prf}).

\begin{table}[t!]
	\centering
	\caption{Eunoia-to-LambdaPi encoding.
		Symbols prefixed with $\msf{eo.}$ are defined in
		our hand-written prelude (\autoref{sec:results});
		$⤳$ and $⤳^{\msf d}$ are the function-space
		symbols introduced in \autoref{sec:lambdapi}.}
	\label{tab:encoding}
	\small
	\begin{tabular}{@{}p{0.40\linewidth}@{\ \ }p{0.52\linewidth}@{}}
		\hline
		\textbf{Eunoia feature} & \textbf{LambdaPi encoding} \\
		\hline
		Type universe $\eotype$
		& term $T : \Set$, carrier $\msf{τ}\,T$ \\
		Function type $\earr{A\ B}$
		& $A ⤳ B$, with $\msf{τ}\,(A ⤳ B) ↪ \msf{τ}\,A → \msf{τ}\,B$ \\
		Dependent type from $\eoquote{x}$
		& $A ⤳^{\msf d} \lam{x}{B}$ (matchable head, unlike $Π$) \\
		Parameter attribute $\attr{implicit}$
		& implicit binder $\imp{x : \msf{τ}\,T}$ \\
		Macro $\kw{define}$
		& $\kw{symbol}\ s\,\vec ρ ≔ t$ \\
		Higher-order application $\eho{t_1}{t_2}$
		& $\lpo{app}\ t_1\ t_2$ \\
		$n$-ary $\eapp{f}{\vec t}$ with attribute
		& nested binary application; nil terminator via rewrite rule \\
		Binder variables ($\eo{var}$)
		& reified into the list type $\lpo{List}$ \\
		Program ($\kw{program}$)
		& $\md{sequential}\ \kw{symbol}$ + one rewrite rule per case (SR-checked) \\
		Inference rule ($\kw{declare-rule}$)
		& $\md{constant}\ \kw{symbol}$ + auxiliary symbol with quote-pattern rewrite when $\attr{args}$ present \\
		Built-in arithmetic
		& prelude rules delegating to $\mathbb Z, \mathbb Q$ \\
		Conclusion guard $\eoreqs{l}{r}{ψ}$
		& $ψ$ applies only when $l ≡ r$ syntactically \\
		Proof command $\kw{assume}$
		& $\md{constant}\ \kw{symbol}$ of type $\lpo{Proof}\,φ$ \\
		Proof command $\kw{step}$
		& defined symbol + $\kw{assert}$ \\
		\hline
	\end{tabular}
\end{table}


% ----------------------------------------------------------
\subsubsection{Translating Commands.}
% ----------------------------------------------------------
Parameter contexts and two auxiliary notations
support the command translation that follows.

\begin{definition}\label{def:ctx}
	Given $\vec ρ = \epar{x_1\ t_1\ \maybe{ν_1}}
		\ldots \epar{x_n\ t_n\ \maybe{ν_n}}$, define:
	$$
		\ctx{\vec ρ}\ ≔\
		ρ_1^★\, \ldots\, ρ_n^★
		\quad\text{where}\quad
		ρ_i^★ =
		\begin{scases}
			\imp{\bar x_i : \El{\tm{t_i}}}
			& \text{if $ν_i = \attr{implicit}$,} \\
			\xpl{\bar x_i : \El{\tm{t_i}}}
			& \text{otherwise.}
		\end{scases}
	$$
\end{definition}
%
We further write $\tm{t}^{\ptrn{}}$ for the result of
applying $\tm{⋅}$ to $t$ and replacing each free
variable $x$ with the pattern variable $\ptrn{x}$;
and $\mbf{impl}(\vec ρ, t)$ for the sublist of
$\ctx{\vec ρ}$ containing each $x_i ∈ \vec ρ$ that
occurs free in $t$, marked as implicit.
%
\begin{definition}\label{def:cmd}
	%
	Let $δ$ be a standard Eunoia command. Then $\cmd{δ}$
	is the list of LambdaPi commands described as follows:
	%
	\begin{enumerate}\itemsep3pt
		\item If $δ$ is a constant declaration,
		      $δ ⊢ s(\vec ρ) : t$, the head of
		      $\cmd δ$ is
		      $\kw{symbol}\ \bar s\ \ctx{\vec ρ}
			      : \El{\tm t}\mtt{;}$.
		      If $δ ⊢ s ∷ α$ and $α$ is
		      $\rcn{t_\nil}$ or $\lcn{t_\nil}$, the
		      following rewrite rule is also produced,
		      where $\lpo{nil}$ returns the nil-terminator
		      of an associative symbol and $\lpo{as}$
		      performs type ascription:
		      $$
			      \kw{rule}\ \prn{\lpo{nil}\ \bar s\ \ptrn T}
			      ↪ \prn{\lpo{as}\ \tm{t_\nil}\ \ptrn T}
		      $$

		\item If $δ$ is a macro definition,
		      $δ ⊢ s(\vec ρ) ≔ t$, then
		      $\cmd{δ}$ is
		      $\kw{symbol}\ \bar s\ \ctx{\vec ρ} :
			      \El{\tm{t'}} ≔ \tm{t}$
		      (with the typing omitted if no $t'$
		      is supplied).

		\item If $δ$ is a program declaration with
		      $δ ⊢ s(\vec ρ) : \earr{\plur t n\, t'}$
		      and cases $\ecase{l_1}{r_1}\ldots\ecase{l_m}{r_m}$,
		      then $\cmd{δ}$ declares $\bar s$,
		      where $T ≔ \El{\tm{ \earr{\plur t n\, t'}}}$:
		      %
		      $$\begin{array}[t]{l}
				      \md{sequential}\ \kw{symbol}\ \bar s\
				      \vec ρ^★ : T
				      \mtt{;}
				      \\[1mm]
				      \begin{array}[t]{@{}l@{\ }l@{\ {↪}\ }l}
					      \kw{rule} & \tm{l_1}^{\ptrn{}} & \tm{r_1}^{\ptrn{}}
					      \\ \cdots \\
					      \kw{with} & \tm{l_m}^{\ptrn{}} & \tm{r_m}^{\ptrn{}}\mtt{;}
				      \end{array}
				      % \quad (j = 1 \ldots m)
			      \end{array}$$
		      %
		      where $\vec ρ^★ ≔ \mbf{impl}(\vec ρ, T)$ binds
		      all of the free parameters in $T$ as implicits.

		\item If $δ$ is a rule declaration with
		      $\attr{args}\ \epar{\plur t n}$ of type
		      $\plur τ n$, then
		      $δ ⊢ s(\vec ρ) :
			      \earr{
				      \eoquote{t_1}\,\ldots\,\eoquote{t_n}
				      \ τ^★}$,
		      where $τ^★$ encodes the premises and
		      conclusion (\autoref{sec:eo-sig}).
		      An auxiliary symbol $s^★$ eliminates the
		      $\lpo{quote}$ types by rewriting to
		      $\tm{τ^★}$ on terms matching the argument
		      patterns; $s$ then takes the arguments
		      explicitly and uses $s^★$ for its return type:
		      $$
			      \begin{array}[t]{l}
				      \kw{symbol}\ \bar{s}^★\
				      \vec σ : \tm{τ_1} → \ldots → \tm{τ_n} → \TYPE \mtt{;}
				      \\[1mm]
				      \kw{rule}\
				      \bar s^★\,\tm{t_1}^{\ptrn{}}\,\ldots\,\tm{t_n}^{\ptrn{}}
				      ↪ \tm{τ^★}^{\ptrn{}}\mtt{;}
				      \\[1mm]
				      \md{constant}\ \kw{symbol}\ \bar s\
				      (α_1 : \tm{τ_1}) \ldots (α_n : \tm{τ_n}) :
				      \prn{\app{\app{{\app{\bar s^★}{α_1}}}{\ldots}}{α_n}}
			      \end{array}
		      $$
	\end{enumerate}
\end{definition}

\begin{example}\label{ex:cmd}
	We illustrate $\cmd{δ}$ on the four declarations
	from \autoref{ex:eo-rules}:
	%
	\begin{lpbox}
		\lstinputlisting[style=lp, aboveskip=0pt, belowskip=0pt]{listings/cmd-constants.lp}%
		\listsep
		\lstinputlisting[style=lp, aboveskip=0pt, belowskip=0pt]{listings/cmd-modus-ponens.lp}%
		\listsep
		\lstinputlisting[style=lp, aboveskip=0pt, belowskip=0pt]{listings/cmd-rule-aux.lp}%
	\end{lpbox}
	%
	\noindent
	The polymorphic equality $\ev{=}$ and the
	associative $\ev{or}$ translate by case~1, the latter
	also producing a nil-terminator rewrite rule.
	$\ev{modus\_ponens}$ has no $\attr{args}$, so case~1
	yields its derived type directly.
	$\ev{cnf\_implies\_pos}$ has the compound argument
	$\attr{args}\epar{\eapp{\ev{=>}}{\ev{F1}\ \ev{F2}}}$,
	triggering case~4: an auxiliary symbol matches the
	argument pattern and rewrites to the translated
	conclusion under $\lpo{Proof}$.
\end{example}

% ----------------------------------------------------------
\subsubsection{Translating Proofs.}
% ----------------------------------------------------------

\begin{definition}\label{def:cmd-prf}
	%
	Let $π$ be a proof command. Then $\cmd{π}$
	is the list of LambdaPi commands such that:
	%
	\begin{enumerate}\itemsep6pt
		\item If $π$ is an assumption,
		      then $π ⊢ s : \eapp{\mtt{Proof}}{φ}$:
		      %
		      $$\md{constant}\ \kw{symbol}\ \bar s :
			      \lpo{Proof}\ \tm{φ}\mtt{;}$$

		\item If $π$ is a step,
		      then $π ⊢ s : \eapp{\mtt{Proof}}{φ}$
		      and $π ⊢ s ≔ \eapp{r}{\plur t m\ \plur p n}$:
		      %
		      \begin{align*}
			       & \kw{symbol}\ \bar s
			      ≔ \bar r\ \tm{t_1}\ \ldots\ \tm{t_m}\
			      \bar p_1\ \ldots\ \bar p_n\mtt{;}
			      \\
			       & \kw{assert}\ ⊢\ \bar s :
			      \lpo{Proof}\ \tm{φ}\mtt{;}
		      \end{align*}
		      %
		      The body applies rule $\bar r$ to its
		      arguments then premises.
		      %
		      The type is checked separately via
		      $\kw{assert}$, which uses LambdaPi's
		      type inference to fill implicit arguments.
	\end{enumerate}
\end{definition}

\begin{example}\label{ex:proof}
	A \texttt{QF\_UF} proof script (top) deriving
	$\false$ from $\eapp{\ev =}{\ev{x0}\ \ev{x0}}$
	via $\ev{eq{-}refl}$ and $\ev{eq\_resolve}$ (CPC
	uses both \texttt{-} and \texttt{\_} in rule names),
	and its LambdaPi translation (bottom) generated by
	\autoref{def:cmd} and \autoref{def:cmd-prf}:
	%
	\begin{eobox}
		\lstinputlisting[style=eo, aboveskip=0pt, belowskip=0pt]{listings/proof.eo}%
		\listsep
		\lstinputlisting[style=lp, aboveskip=0pt, belowskip=0pt]{listings/proof.lp}%
	\end{eobox}
	%
	The preamble's $\ev{U}$, $\ev{x0}$ translate by
	case~1 and the $\kw{define}$s by case~2 of
	\autoref{def:cmd}; the assumption $\ev{@p1}$
	becomes a $\md{constant}$ symbol of type
	$\lpo{Proof}\ \tm{φ}$ (case~1 of
	\autoref{def:cmd-prf}); each step applies its
	rule to arguments then premises, with $\kw{assert}$
	verifying the conclusion via type inference.
\end{example}

\end{document}
